% k0x_fill_Kh2: the two branches of the proof of hyp:k02:Kh and its assembly (input from sec/k0x_fill_Kh.tex).

\begin{proposition}[the branch $R\le p$: zone~I]\label{prop:k02:KhI}
Let $H$ be a member of $\mathcal N_h$ and $\tau$ a localisation of $H$ in zone~I; by \cref{lem:a02:Z}(i) this covers every localisation with $R\le p$. Then, for each model, on $\supp\chi_j$,
\begin{equation}\label{eq:k02:I1}
|E^{(j)}(H;\tau)|\ \le\ C(1+k-j)^{p_0}\,\eta^{\rho(H)}\,p^{n_1}\,e^{-\delta R},
\end{equation}
which is \cite[(3.39)]{K1}, and across the diagonal pair, on the common unit lattice,
\begin{equation}\label{eq:k02:I2}
|E'^{(j+1)}(H;\tau)-E^{(j)}(H;\tau)|\ \le\ C\,\eta^{\rho(H)-1/2}\,p^{n_1}\big[5^{-\gamma j}+\eta^{3+\sigma+1/2}\big]\,e^{-\delta R}.
\end{equation}
The constants depend on the point of King's range, $J$, $\bar n$, $\norm h_{C^1}$ and, for members with curvature vertices, $\max\norm{f_{\mu\nu}}_{C^{2,\alpha}}$; not on $j$, $k$, $\tau$.
\end{proposition}
\begin{proof}
By \cref{lem:a02:I}, which uses only that the localisation is in zone~I, $E(H)=E(H^{\mathrm{free}})+\mathcal R$ with $H^{\mathrm{free}}$ the same graph on the free $f_D$ flow, every term of $\mathcal R$ carrying an image line of any prescribed degree with the factor $e^{-\delta_0p/2}\le C\eta^{3+\sigma+1/2}$, and the free-boundary counterterm renormalising $H^{\mathrm{free}}$ exactly. For $H^{\mathrm{free}}$ we run King's proof of \cite[Thm.~3.5, Prop.~3.6]{K1} as in the proof of \cref{thm:a02:B4}, with the following points for the $h$-vertex. \emph{Census:} every connected component of an ordered prefix that contains the $h$-vertex has positive degree, the vertex \cite[(3.21)]{K1} carrying the coupling power $2$ against at most two fields, except the field-independent $h$-linear graphs of the three lowest orders $P_0$ \cite[Thm.~3.5]{K1}, which form the family $\mathcal T_h$ of \cref{app:a04:def:families} and are not in $\mathcal N_h$; the $h$-free components obey the degree count of \cref{thm:a02:T}, with the Ward rows by \cref{thm:a02:WR}, and the components with block points are convergent by \cref{lem:a02:BLK}. \emph{Bound:} the ordered multiscale bound \cref{thm:a02:S3}(a) with the root held at the $h$-vertex gives \eqref{eq:k02:I1}; the power of $\eta$ is $\rho(H)$ by \cite[p.~665]{K1}, the fields on $\supp\chi_j$ being bounded by \cite[(3.2)]{K1}, which is the content of \eqref{eq:k02:margin}, and the logarithm comes from the infrared sums as in \cref{thm:a02:B4}. \emph{Comparison:} \cref{thm:a02:S3}(b), one factor at a time \cite[pp.~663--665]{K1}: the lines by \cref{prop:a02:FL}, the $h$-vertex by the sampling of the same smooth $h$ at points of the two frames at distance $\le\sqrt3\,\eps_k$, whose difference is $\le C\eps_k\norm h_{C^1}\le C5^{-\gamma j}$, and the finer slices through the finest index sum; each error term is the graph with one element weakened, which costs the factor $\eta^{-1/2}$ in \eqref{eq:k02:I2}. The identities \cite[(3.43)--(3.55)]{K1} act on the $h$-vertex only through the phases at its scalar ends, which the gauge transformation \cite[(3.46)]{K1} carries. Members with curvature vertices: \cref{lem:a02:SBP}(b)--(d) as in \cref{thm:a02:B4}. The $h$-vertex is bounded by $\norm h_\infty$, and only through it does $h$ enter the constants of \eqref{eq:k02:I1}.
\end{proof}

\begin{proposition}[the branch $R>p$: the cut by clusters]\label{prop:k02:KhR}
Let $\eta\le\eta_2$, let $H$ be a member of $\mathcal N_h$ and $\tau$ a tuple of $H$ with at least two clusters. By \cref{lem:k02:Zh}(c) this covers every zone-B tuple, and by \cref{lem:a02:CL}(c) and \eqref{eq:k02:range} every tuple with $R>p$. Then, for each model separately, on $\supp\chi_j$,
\begin{equation}\label{eq:k02:R1}
|E^{(j)}_N(H;\tau)|\ \le\ C(1+k-j)^{p_0}\,\eta^{\rho(H)}\,p^{n_1}\,e^{-\delta\,\dist(P(\tau))},
\end{equation}
$E_N$ the field-dependent part (the whole value for a field-independent member); this is \cite[(3.38)]{K1} at the tuple. Moreover, for every unit cube $y_*$ and $0<\delta'\le\delta/2$,
\begin{equation}\label{eq:k02:R2}
\sum_{\substack{\tau\ \text{in zone B}\\ \text{with an $h$-vertex in }y_*}}|E^{(j)}_N(H;\tau)|\,e^{\delta'R(\tau)}\ \le\ C(1+k-j)^{p_0}\,\eta^{\rho(H)}\,p^{n_1}\,e^{-\delta/(8\eta)} .
\end{equation}
The constants depend on the point of King's range, $J$, $\bar n$, $\norm h_\infty$ and, for members with curvature vertices, $\max\norm{f_{\mu\nu}}_{C^{2,\alpha}}$.
\end{proposition}
\begin{proof}
At a tuple with at least two clusters every connected component of every cluster graph carries a cut-line leg (\cref{cor:a02:DEGall}(c)), so the vacuum members and the component-wise cancellation of \cite[p.~667]{K1} do not arise, and $E(H;A^{(j)};\tau)$ and $E(H;0;\tau)$ are bounded separately; at zero background the cut lines are the cube propagators at $A=0$ and $B_K=0$. Fix one of the two terms.

\emph{Step 0: regularised cut lines.} Replace every cut line $\ell$ by $\ell^\flat$; by \cref{lem:a02:XL}(a) this changes nothing.

\emph{Step 1: reduction per cluster.} Each $H_K$ is reduced by \cite[(3.43)--(3.55)]{K1} with its own box $\Omega_K$, centre $x_0^K$, gauge $B_K$ and field $\tilde A_K=A^{(j)}-B_K$ on $\Omega_K$ (\cref{lem:a02:CL}), its cut lines keeping the background and receiving the unitary factors \cite[(3.46)]{K1} at their ends. The expansion of the reduced propagator in $\tilde A_K$ converges on $\Omega_K$: $e_j\sup_{\Omega_K}|\tilde A_K|\le Ce_jp\operatorname{diam}\Omega_K\le C_\flat e\eta^{1/2}p^2\le\frac12$ by the difference bound \cite[(3.45)]{K1}, \cref{lem:k02:Zh}(e) and \eqref{eq:k02:range}: this is the reduction of the branch $R\le p$ on a box of side $O(p)$, run inside each cluster. By \cref{lem:a02:CL}(a) the partition, the boxes, the gauges and the cut status of every line are the same for every member of every family met in the nesting and for every position, so renormalisation families are formed inside each $H_K$. A zone-I cluster, in particular every cluster containing a source vertex (\cref{lem:k02:Zh}(d)), satisfies $\dist(\square_K,\partial\Omega_K)\ge2p$ and $\Omega_K\cap F=\emptyset$, so \cref{lem:a02:I} applies to $H_K$: it is its free $f_D$-flow version plus terms carrying an image line of any prescribed degree, with the cut lines as external legs; the vertices of the expansion in $\tilde A_K$ that lie in $\square'_K\setminus\square_K$ are joined to a localised vertex of $K$ by a path of lines one of which is long, and that line supplies a credit of any order, as in step (2) of the proof of \cref{thm:a02:WR} run inside $H_K$ (\cref{thm:a02:WRall}). A zone-B cluster is source-free and is rewritten by \cref{def:a02:RC} within $\Omega_K$.

\emph{Step 2: the cut-line factors.} Every evaluation point of a cut-line end in an output term is joined to a localised anchor of its own cluster by at most $C(\bar n)$ lines of that cluster's rewritten units (\cref{thm:a02:WRall}), and by \cref{lem:a02:XL}(b) the phase-dressed cut line satisfies $|O_uO_v\ell^{\flat,\mathrm{ph}}|\le Ce^{-\delta_2d(K,K')}\times$(transfer factors on the lines of those paths), for all operations $O_u,O_v$ of order $\le1+\alpha$ at its ends, and $\le C$ crudely. The quotient graph (clusters as nodes, cut lines as edges) is connected because $H$ is (\cref{lem:a02:CL}(d)); bound the cut lines of one spanning tree $T$ of it by the first bound and the others by $C$.

\emph{Step 3: factorisation.} The output term is bounded by $\prod_K\mathfrak E_K\times\prod_{\ell\in T}Ce^{-\delta_2d(K,K')}\times C^{\#\text{cut lines}}$, where $\mathfrak E_K$ depends on the positions of $K$ only; the transfer factors cost one extra decay draw per line of a cluster, within the budget of \cref{thm:a02:WRall}.

\emph{Step 4: the cluster factors.} \emph{Zone-I clusters.} The cluster graph is a free $f_D$-flow graph with its cut lines as undifferentiated external legs of unit scale (\cref{lem:a02:XL}(b)). For the cluster containing the $h$-vertex, the census is that of \cref{prop:k02:KhI}; since the cluster carries a cut leg, every component has an external leg and the excepted field-independent family does not occur. For the source-free zone-I clusters, the degree count holds at every tuple by \cref{thm:a02:T} and \cref{thm:a02:WRall}, and block points by \cref{lem:a02:BLK}. The multiscale bound \cref{thm:a02:S3}(a) with the root held, summed over the root in $\Omega_K$ (volume $\le Cp^3$ by \cref{lem:a02:CL}(c)), gives $\sum_{\text{positions of }K}\mathfrak E_K\le C\eta^{\rho_K}p^{n_K}(1+k-j)^{p_{0,K}}$, with $\rho_K$ the coupling power of the vertices of $K$. \emph{Zone-B clusters.} The cluster is source-free (\cref{lem:k02:Zh}(d)); apply \cref{thm:a02:Bsc}(a) and \cref{thm:a02:BVP} with no vertex held to the renormalised family of $H_K$, localised at $P_K$ with the box $\Omega_K$ of side $O(p)$; the hypothesis ``at least one external field'' is the cut leg. The proof of \cref{thm:a02:Bsc} uses an external scalar leg only in the following ways, each of which a cut line satisfies: through its sup and $(1,\alpha)$-norms on $\Omega_K$ up to the faces, in the Taylor and transport steps, the H\"older splits and the ordered induction (\cref{lem:a02:XL}(b)); and, when the charging rule of \cref{lem:a02:CH} charges the leg at a $D'=-1$ remainder vertex $v$, through the vanishing factor $m_\alpha(v;1)$ of \cref{lem:a02:H}(ii), which a cut line supplies by \cref{lem:k02:face}(a). A cut line joining two zone-B clusters may be charged at both ends, by the charging rules of the two clusters; then \cref{lem:k02:face}(b) supplies both factors. After the borrowing of \cref{lem:a02:CH} the cut-line factor left in the output term is $\ell/[m_\beta(u';1)m_\beta(v';1)]$ (or with one factor), an undifferentiated evaluation bounded by $C\min(1,t(u'))^{1-\beta}\min(1,t(v'))^{1-\beta}e^{-\delta\dist}\le Ce^{-\delta\dist}$ by \cref{lem:k02:face}, which Steps~2--3 treat like every other cut-line factor. The resulting bound is again $\le C\eta^{\rho_K}p^{n_K}(1+k-j)^{p_{0,K}}$, the logarithm coming from the infrared sums inside $\Omega_K$.

\emph{Step 5: the tree.} By \cref{lem:a02:CL}(d), $\prod_{\ell\in T}e^{-\delta_2d(K,K')}\le e^{\delta_2c_0N}e^{-\delta_2\dist(P(\tau))}$.

\emph{Step 6: exponents.} $\rho(H)=\sum_K\rho_K$, the count being vertex-local (an unlocalised mass-counterterm vertex on a cut dressed line keeps its own coupling power); $n_1\defeq\sum_Kn_K$ and $p_0\defeq\sum_Kp_{0,K}$ plus the dressed-line factor of \cref{lem:a02:XL}(b); the number of clusters is $\le N_*$ and $\delta\defeq\delta_2$. This is \eqref{eq:k02:R1}.

\emph{The root sum.} On the tuples of \eqref{eq:k02:R2}, $R(\tau)\ge\eta^{-1}/2$ by \cref{lem:k02:Zh}(b). The number of tuples of $H$ with a point in $y_*$ and $R\in[n,n+1)$ is $\le C(n+1)^{3N_*}$, and $e^{\delta'R}e^{-\delta R}\le e^{-\delta R/2}$. So the sum is $\le C(1+k-j)^{p_0}\eta^{\rho(H)}p^{n_1}\sum_{n\ge\eta^{-1}/2-1}(n+1)^{3N_*}e^{-\delta n/2}\le C'(1+k-j)^{p_0}\eta^{\rho(H)}p^{n_1}e^{-\delta/(8\eta)}$, since $(n+1)^{3N_*}e^{-\delta n/4}\le C$ and $\eta^{-1}\ge8$. Members with curvature vertices: the curvature vertex lies in a zone-I cluster (\cref{lem:k02:Zh}(d)), is invariant under the reduction (\cref{lem:a02:SBP}(b)) and is bounded through \cref{lem:a02:SBP}(c).
\end{proof}

\begin{proof}[Proof of \cref{hyp:k02:Kh}]
Let $H$ be a member of $\mathcal N_h$ and $J_*\le j\le k-J_*$, so that $\eta\le\eta_2$. \emph{First clause.} \eqref{eq:k02:I1} is \cite[(3.39)]{K1} for the localisations in zone~I, which are those at distance from the faces larger than the small-field scale. \emph{Second clause.} A tuple with one cluster has $R\le c_0(N_*-1)^2\le p$, lies in zone~I by \cref{lem:a02:Z}(i), and obeys \eqref{eq:k02:I1}; a tuple with at least two clusters obeys \eqref{eq:k02:R1}. Summed over the localisations sharing the cube $y_*$ of an $h$-vertex: the zone-I tuples by $\sum_n(n+1)^{3N_*}e^{-(\delta-\delta')n}\le C$, the zone-B tuples by \eqref{eq:k02:R2}. This is \cite[(3.38)]{K1} with the factor $(1+k-j)^{p_0}$. \emph{Third clause.} Split the localisations at $R=p$. Those with $R\le p$ are in zone~I; sum \eqref{eq:k02:I2} by \cref{lem:a02:SC} (branch $R\le p$): the result is $\le C\eta^{\rho(H)-1/2}p^{n_1}(1+k)^{p_0}\eta^{-3}|\CJ|[5^{-\gamma j}+\eta^{3+\sigma+1/2}]$, whose first term is $\le C5^{-\gamma'j}\eta^{-\beta'}|\CJ|$ for every $\gamma'<\gamma$ and $\beta'>\frac72-\rho(H)$, the powers $p^{n_1}(1+k)^{p_0}$ being absorbed on the diagonal sequence as in \cref{lem:a02:SC}, and whose second term is $\le C\eta^{\rho(H)+\sigma}p^{n_1}(1+k)^{p_0}|\CJ|\le C\eta^{\sigma'}|\CJ|$ for every $\sigma'<\sigma$. Those with $R>p$ have at least two clusters; bound the two members of the pair separately by \eqref{eq:k02:R1}, with the roots at the $h$-vertex, $\le C\eta^{-3}|\CJ|$ of them. For the zone-I tuples among them the factor $e^{-(\delta-\delta')p}$ of \cref{lem:a02:SC} (branch $R>p$) is $\le C\eta^{3+\sigma}$, since $p(\eta)$ grows faster than every logarithm, and the sum is $\le C\eta^{-3}|\CJ|\eta^{\rho(H)}p^{n_1}(1+k-j)^{p_0}\eta^{3+\sigma}\le C\eta^{\sigma'}|\CJ|$; for the zone-B tuples \eqref{eq:k02:R2} gives $\le C\eta^{-3}|\CJ|\eta^{\rho(H)}p^{n_1}(1+k-j)^{p_0}e^{-\delta/(8\eta)}\le C_N\eta^N|\CJ|$ for every $N$. In both cases $\rho(H)\ge0$, and the external fields on $\supp\chi_j$ are accounted for in $\rho(H)$ by \eqref{eq:k02:margin}. Adding the three parts gives the third clause with $C$ independent of $j$ and $k$. The constants depend on the point of King's range, so the lemma holds at every point of the range, in particular at the second mass $M_2$ of \cref{ssec:k02:setting}.
\end{proof}
